\documentclass[11pt]{article}
\usepackage[margin=1in]{geometry}
\usepackage{amsmath, amssymb, amsthm}
\usepackage{enumitem}
\usepackage{booktabs}
\usepackage{microtype}
\usepackage{tikz}
\usetikzlibrary{calc}
\usepackage{hyperref}

\newcommand{\X}[2]{X_{#1#2}}
\newcommand{\Z}{\mathcal{Z}}
\newtheorem{theorem}{Theorem}
\newtheorem{lemma}{Lemma}
\newtheorem{proposition}{Proposition}
\theoremstyle{definition}
\newtheorem{definition}{Definition}
\newtheorem*{remark}{Remark}

\title{Locality emerges from hidden zeros in \texorpdfstring{$\mathrm{tr}(\phi^3)$}{tr(phi^3)}:\\
a by-hand proof at four, five, and six points}
\author{Jony}
\date{\today}

\begin{document}
\maketitle

\begin{abstract}
We prove that the cyclic hidden one-zeros of color-ordered $\mathrm{tr}(\phi^3)$
tree amplitudes force every non-local coefficient of the general rational
ansatz to vanish, at $n=4, 5, 6$ external legs. The proof is entirely by hand:
we use only the kinematic mesh identity, geometric series Laurent expansions,
and the fact that rational functions in independent variables vanish iff
each coefficient on a basis of monomials vanishes.
The mechanism is geometric: each cyclic zero kills exactly the non-local
terms whose chord factors involve the ``enclosed vertex'' of a designated
short diagonal (the \emph{special variable}). Every non-local coefficient is
killed by at least one zero --- and as it happens, by exactly two.
\end{abstract}

\tableofcontents

%================================================================
\section{Setup}

\subsection{Chords, the ansatz, and locality}

Label the external legs cyclically by $1,\ldots,n$. The planar Mandelstam
variables are
\[
\X{i}{j} \;=\; (p_i+p_{i+1}+\cdots+p_{j-1})^2,\qquad 1\le i<j\le n,\ j-i\ge 2,\ (i,j)\ne(1,n).
\]
Geometrically, $\X{i}{j}$ is the chord from vertex $i$ to vertex $j$ of the
convex $n$-gon. There are $n(n-3)/2$ such chords ($2,5,9$ for $n=4,5,6$).

A \emph{triangulation} of the $n$-gon is a maximal set of non-crossing chords;
it has $n-3$ chords. The tree amplitude is the sum over triangulations of the
product of chord propagators,
\[
A_n \;=\;\sum_{T\text{ triangulation}}\ \prod_{(i,j)\in T}\frac{1}{\X{i}{j}}.
\]
The number of triangulations is the Catalan number $C_{n-2}$ ($2, 5, 14$
for $n=4,5,6$).

We consider the most general rational function $B$ with the correct dimension
($-2(n-3)$) and only planar simple poles:
\[
B \;=\; \sum_{\text{multisets }(a_1,\ldots,a_{n-3})}
\frac{a_{a_1\ldots a_{n-3}}}{X_{a_1}\cdots X_{a_{n-3}}}.
\]
The sum is over multisets (unordered with repetition) of size $n-3$ from
the set of chord indices. There are $\binom{d+n-4}{n-3}$ such multisets,
where $d=n(n-3)/2$ is the number of chords. We get $2, 15, 165$ multisets
for $n=4,5,6$.

A multiset is \emph{local} if its chords form a triangulation; the remaining
multisets are \emph{non-local} (either repeated chords or two chords that
cross). Locality is the statement that all non-local coefficients vanish.

\subsection{The kinematic mesh identity}

The non-planar two-particle invariants $c_{a,b}=2p_a\cdot p_b$ ($a\ne b$,
non-adjacent) are not independent of the planar variables. They satisfy
\[
\boxed{\ c_{a,b} \;=\; \X{a}{b+1} + \X{a+1}{b} - \X{a}{b} - \X{a+1}{b+1}\ }
\]
(adjacent chords like $\X{a}{a+1}$ are zero by masslessness $p_a^2=0$).
Setting $c_{a,b}=0$ thus imposes a linear relation on the planar chords.

\subsection{The hidden one-zeros}

For each leg $r\in\{1,\ldots,n\}$, the \emph{hidden one-zero} $\Z_r$ is
defined by setting all $c$-invariants involving leg $r$ to zero:
\[
\Z_r:\qquad c_{r,r+2}=c_{r,r+3}=\cdots=c_{r,r+n-2}=0\ \pmod n.
\]
There are $n-3$ such constraints. The starting fact (proved elsewhere; here
taken as given) is that the tree amplitude $A_n$ vanishes on each $\Z_r$.
Our problem is the \emph{converse} for the general ansatz: assuming
$B|_{\Z_r}=0$ for all $r$, show that all non-local $a$-coefficients vanish.

\subsection{The geometric heart: enclosed vertices}

Of the two arcs into which a chord $\X{i}{j}$ divides the polygon's
boundary, the \emph{enclosed arc} is the shorter one. The
\emph{enclosed vertex set} is
\[
E(\X{i}{j}) \;=\; \{\text{vertices strictly between $i$ and $j$ on the shorter arc}\}.
\]
We call $|E(\X{i}{j})|$ the \emph{class} of the chord:
class 1 chords (``short diagonals'') enclose $1$ vertex, class 2 (``long
diagonals'') enclose $2$, etc.

\begin{lemma}[Crossing $=$ enclosed-vertex sweep]
\label{lem:crossing}
Two chords $\X{i}{j}$ and $\X{r}{s}$ of a convex polygon cross in the
interior iff exactly one endpoint of $\X{r}{s}$ lies in $E(\X{i}{j})$
(equivalently in the enclosed arc of $\X{i}{j}$), the other in the outer arc.
\end{lemma}

\begin{proof}
Two chords of a convex polygon cross iff their four endpoints alternate around
the boundary. Fix $\X{i}{j}$ and consider the two arcs it cuts off (one
``enclosed,'' one ``outer''). The endpoints of $\X{r}{s}$ either both lie in
the same arc (no alternation, no crossing) or one lies in each (alternation,
crossing).
\end{proof}

\paragraph{Consequence.}
Every non-local pair of chords --- i.e.\ a crossing pair --- has the form
``one chord with enclosed vertex $v$, another chord starting at $v$.'' The
hidden zeros, as we'll see, are tailor-built to kill these.

%================================================================
\section{Four points: trivial warm-up}

The square has just two chords, $\X{1}{3}$ and $\X{2}{4}$. The ansatz
\[
B = \frac{a_{\X{1}{3}}}{\X{1}{3}} + \frac{a_{\X{2}{4}}}{\X{2}{4}}
\]
has $2$ coefficients, both \emph{local} (each chord by itself is a triangulation).
There are no non-local terms to kill. A single zero $\Z_1$ gives
$c_{1,3}=\X{1}{4}+\X{2}{3}-\X{1}{3}-\X{2}{4}=-\X{1}{3}-\X{2}{4}=0$, hence
$\X{2}{4}=-\X{1}{3}$, and $B|_{\Z_1}=(a_{\X{1}{3}}-a_{\X{2}{4}})/\X{1}{3}=0$
forces $a_{\X{1}{3}}=a_{\X{2}{4}}$. This is unitarity (equal residues on both
channels), not locality. We move on to $n=5$, where the question becomes
nontrivial.

%================================================================
\section{Five points: complete proof}
\label{sec:n5}

\subsection{The pentagon's chords and crossings}

The five chords of the pentagon, listed with their enclosed-vertex sets:
\begin{center}
\begin{tabular}{cccccc}
\toprule
Chord & $\X{1}{3}$ & $\X{1}{4}$ & $\X{2}{4}$ & $\X{2}{5}$ & $\X{3}{5}$ \\
$E(\cdot)$ & $\{2\}$ & $\{5\}$ & $\{3\}$ & $\{1\}$ & $\{4\}$ \\
\bottomrule
\end{tabular}
\end{center}

Every chord has class $1$ (encloses one vertex). At $n=5$ there are no
class-2 chords yet; they appear starting at $n=6$.

The five \emph{crossing pairs} (verified by Lemma \ref{lem:crossing}): for each
pair, one chord's enclosed vertex is an endpoint of the other.
\begin{center}
\begin{tabular}{cccc}
\toprule
Pair & $E$-witnesses & Pair & $E$-witnesses \\
\midrule
$\X{1}{3}\times\X{2}{4}$ & $2\in\X{2}{4},\ 3\in\X{1}{3}^c$ &
$\X{1}{4}\times\X{3}{5}$ & $5\in\X{3}{5}^c,\ 4\in\X{1}{4}^c$ \\
$\X{1}{3}\times\X{2}{5}$ & $2\in\X{2}{5},\ 1\in\X{1}{3}^c$ &
$\X{2}{4}\times\X{3}{5}$ & $3\in\X{3}{5}^c,\ 4\in\X{2}{4}^c$ \\
$\X{1}{4}\times\X{2}{5}$ & $5\in\X{2}{5}^c,\ 1\in\X{1}{4}^c$ & & \\
\bottomrule
\end{tabular}
\end{center}

The other $5$ unordered pairs of distinct chords share an endpoint (so
non-crossing) and are exactly the $5$ fan triangulations.

\subsection{The 15-coefficient ansatz}

With $n-3=2$ chord factors per monomial, the ansatz is
\[
B = \sum_{1\le i\le j\le 5}\frac{a_{ij}}{X_i X_j},\quad\text{15 coefficients},
\]
where the indexing runs over multisets of chord-indices. Of the $15$ coefficients:
\begin{itemize}[nosep,leftmargin=2em]
\item $5$ \emph{tree triples}: one for each fan triangulation.
\item $5$ \emph{double poles}: $a_{\X{i}{j}\X{i}{j}}$, one per chord.
\item $5$ \emph{crossing pairs}: one $a_{\X{i}{j}\X{k}{\ell}}$ per crossing.
\end{itemize}
The latter $10$ are the non-local coefficients we must kill.

\subsection{The cyclic hidden one-zeros, derived by hand}

We derive each $\Z_r$ from $c_{r,r+2}=c_{r,r+3}=0$ using the mesh identity.

\paragraph{$\Z_1$:} $c_{1,3}=c_{1,4}=0$.
\begin{align*}
c_{1,3}&=\X{1}{4}+\X{2}{3}-\X{1}{3}-\X{2}{4}=\X{1}{4}-\X{1}{3}-\X{2}{4}=0,\\
c_{1,4}&=\X{1}{5}+\X{2}{4}-\X{1}{4}-\X{2}{5}=\X{2}{4}-\X{1}{4}-\X{2}{5}=0.
\end{align*}
(We used $\X{2}{3}=\X{1}{5}=0$, since $\X{a}{a+1}$ corresponds to the
massless particle squared.) Solving for $\X{2}{4}$ and $\X{2}{5}$ (the two
chords from the enclosed vertex $2$ of $\X{1}{3}$):
\begin{equation}
\Z_1:\qquad\boxed{\X{2}{4}=\X{1}{4}-\X{1}{3},\qquad \X{2}{5}=-\X{1}{3}.}
\end{equation}
The chord $\X{1}{3}$ appears on the right of \emph{both} substitutions: it
is the \textbf{special variable} of $\Z_1$.

\paragraph{$\Z_2$:} $c_{2,4}=c_{2,5}=0$.
$c_{2,4}=\X{2}{5}+\X{3}{4}-\X{2}{4}-\X{3}{5}=\X{2}{5}-\X{2}{4}-\X{3}{5}=0$.
$c_{2,5}=\X{2}{6}+\X{3}{5}-\X{2}{5}-\X{3}{6}$, where $\X{2}{6}=\X{2}{1}=0$
and $\X{3}{6}=\X{3}{1}=\X{1}{3}$, so $c_{2,5}=\X{3}{5}-\X{2}{5}-\X{1}{3}=0$.
Solving:
\begin{equation}
\Z_2:\qquad\boxed{\X{3}{5}=\X{2}{5}-\X{2}{4},\qquad \X{1}{3}=-\X{2}{4}}, \quad\text{special }=\X{2}{4}.
\end{equation}

\paragraph{$\Z_3$:} $c_{3,5}=c_{3,1}=0$ (cyclic).
$c_{3,5}=\X{3}{6}+\X{4}{5}-\X{3}{5}-\X{4}{6}=\X{1}{3}-\X{3}{5}-\X{1}{4}=0$, so
$\X{1}{4}=\X{1}{3}-\X{3}{5}$.
$c_{3,1}=\X{3}{2}+\X{4}{1}-\X{3}{1}-\X{4}{2}=\X{1}{4}-\X{1}{3}-\X{2}{4}=0$, so
$\X{2}{4}=\X{1}{4}-\X{1}{3}=-\X{3}{5}$.
\begin{equation}
\Z_3:\qquad\boxed{\X{1}{4}=\X{1}{3}-\X{3}{5},\qquad \X{2}{4}=-\X{3}{5}}, \quad\text{special }=\X{3}{5}.
\end{equation}

\paragraph{$\Z_4$:} $c_{4,1}=c_{4,2}=0$.
$c_{4,1}=\X{4}{2}+\X{5}{1}-\X{4}{1}-\X{5}{2}=\X{2}{4}-\X{1}{4}-\X{2}{5}=0$, so $\X{2}{5}=\X{2}{4}-\X{1}{4}$.
$c_{4,2}=\X{4}{3}+\X{5}{2}-\X{4}{2}-\X{5}{3}=\X{2}{5}-\X{2}{4}-\X{3}{5}=0$, so $\X{3}{5}=\X{2}{5}-\X{2}{4}=-\X{1}{4}$.
\begin{equation}
\Z_4:\qquad\boxed{\X{2}{5}=\X{2}{4}-\X{1}{4},\qquad \X{3}{5}=-\X{1}{4}}, \quad\text{special }=\X{1}{4}.
\end{equation}

\paragraph{$\Z_5$:} $c_{5,2}=c_{5,3}=0$.
$c_{5,2}=\X{5}{3}+\X{1}{2}-\X{5}{2}-\X{1}{3}=\X{3}{5}-\X{2}{5}-\X{1}{3}=0$, so $\X{1}{3}=\X{3}{5}-\X{2}{5}$.
$c_{5,3}=\X{5}{4}+\X{1}{3}-\X{5}{3}-\X{1}{4}=\X{1}{3}-\X{3}{5}-\X{1}{4}=0$, so $\X{1}{4}=\X{1}{3}-\X{3}{5}=-\X{2}{5}$.
\begin{equation}
\Z_5:\qquad\boxed{\X{1}{3}=\X{3}{5}-\X{2}{5},\qquad \X{1}{4}=-\X{2}{5}}, \quad\text{special }=\X{2}{5}.
\end{equation}

\subsection{The pattern: every zero is a rotation of the previous}
\label{sec:n5-pattern}

Every zero has the same structure. It will help to name the pieces precisely;
this vocabulary will carry over unchanged to $n=6$ and beyond.

\begin{definition}[Canonical vocabulary for $\Z_r$]
\label{def:vocabulary}
Looking at $\Z_r$ in its solved (canonical) form, we call:
\begin{itemize}[nosep,leftmargin=2em]
\item The \textbf{special variable} $X_*$: the chord appearing on the
right-hand side of \emph{every} substitution. At $\Z_r$ it is always the
short diagonal $\X{r}{r+2}$, whose enclosed vertex is $r+1$.
\item The \textbf{substituted chords} $S_r$: the $n-3$ chords on the
left-hand sides of the substitutions. These are exactly the chords starting
at the enclosed vertex $r+1$ of $X_*$.
\item A substituted chord $X_s$ is \textbf{non-bare} if its substitution
has the form $X_s = X_f - X_*$ with $X_f$ a genuinely free chord.
The chord $X_f$ is called the \textbf{companion} of $X_s$, denoted $f_s$.
\item A substituted chord is \textbf{bare} if its substitution is $X_s=-X_*$
(no free chord on the RHS; companion is zero).
\item The \textbf{free set} $F_r$: the $d-(n-3)-1$ chords that are neither
substituted nor special.
\end{itemize}
\end{definition}

At $n=5$ with $|F_r|=2$, each zone has exactly $1$ non-bare substitution
with its $1$ companion in $F_r$, and $1$ bare substitution. Tabulating:

\begin{center}
\renewcommand{\arraystretch}{1.2}
\begin{tabular}{c|c|c|c|c}
\toprule
Zone & Special $X_*$ & Non-bare $X_s$ & Companion $f_s$ & Bare \\
\midrule
$\Z_1$ & $\X{1}{3}$ & $\X{2}{4}$ & $\X{1}{4}$ & $\X{2}{5}$ \\
$\Z_2$ & $\X{2}{4}$ & $\X{3}{5}$ & $\X{2}{5}$ & $\X{1}{3}$ \\
$\Z_3$ & $\X{3}{5}$ & $\X{1}{4}$ & $\X{1}{3}$ & $\X{2}{4}$ \\
$\Z_4$ & $\X{1}{4}$ & $\X{2}{5}$ & $\X{2}{4}$ & $\X{3}{5}$ \\
$\Z_5$ & $\X{2}{5}$ & $\X{1}{3}$ & $\X{3}{5}$ & $\X{1}{4}$ \\
\bottomrule
\end{tabular}
\end{center}

\subsection{The kill technique, in full detail at \texorpdfstring{$\Z_1$}{Z1}}

We work out $\Z_1$ in full and verify the kills by hand. The other four
zeros follow by cyclic rotation; their kills can be read off.

For $\Z_1$, the special is $X_*=\X{1}{3}$, the substitutions are
$\X{2}{4}=\X{1}{4}-\X{1}{3}$ (non-bare, companion $\X{1}{4}$) and
$\X{2}{5}=-\X{1}{3}$ (bare). The free set is $F_1=\{\X{1}{4},\X{3}{5}\}$.

After substitution, $B|_{\Z_1}$ is a rational function of three variables:
$\X{1}{3},\X{1}{4},\X{3}{5}$. The condition $B|_{\Z_1}=0$ means this
rational function vanishes identically. We extract single-coefficient
kills by Laurent-expanding in the special variable $\X{1}{3}$, treating
$\X{1}{4},\X{3}{5}$ as fixed.

\paragraph{Step 1: Order $\X{1}{3}^0$ (the limit $\X{1}{3}\to\infty$).}

As $\X{1}{3}\to\infty$ with $\X{1}{4},\X{3}{5}$ held fixed:
\[
\X{2}{4}=\X{1}{4}-\X{1}{3}\to-\infty,\quad \X{2}{5}=-\X{1}{3}\to-\infty,
\]
so $1/\X{2}{4}\to 0$ and $1/\X{2}{5}\to 0$. Thus every monomial of $B$
containing any factor of $\X{1}{3}$, $\X{2}{4}$, or $\X{2}{5}$ vanishes in
the limit. The surviving monomials use only the two free chords $\X{1}{4},\X{3}{5}$:
\[
\lim_{\X{1}{3}\to\infty} B|_{\Z_1}
\;=\;
\frac{a_{\X{1}{4}\X{1}{4}}}{\X{1}{4}^2} +
\frac{a_{\X{3}{5}\X{3}{5}}}{\X{3}{5}^2} +
\frac{a_{\X{1}{4}\X{3}{5}}}{\X{1}{4}\,\X{3}{5}}.
\]
This must vanish identically as a function of two independent variables
$\X{1}{4},\X{3}{5}$. The three rational functions $1/\X{1}{4}^2$, $1/\X{3}{5}^2$,
$1/(\X{1}{4}\X{3}{5})$ are linearly independent. Hence each coefficient must be zero:
\[
\boxed{\quad a_{\X{1}{4}\X{1}{4}}=a_{\X{3}{5}\X{3}{5}}=a_{\X{1}{4}\X{3}{5}}=0.\quad}
\quad\text{($\Z_1$, three order-0 kills)}
\]

\paragraph{Step 2: Order $\X{1}{3}^{-1}$ (coefficient of $1/\X{1}{3}$).}

Expand the substituted chords as geometric series:
\[
\frac{1}{\X{2}{4}}
=\frac{1}{\X{1}{4}-\X{1}{3}}
=-\frac{1}{\X{1}{3}}\cdot\frac{1}{1-\X{1}{4}/\X{1}{3}}
=-\frac{1}{\X{1}{3}}-\frac{\X{1}{4}}{\X{1}{3}^2}-\frac{\X{1}{4}^2}{\X{1}{3}^3}-\cdots,
\]
\[
\frac{1}{\X{2}{5}}=-\frac{1}{\X{1}{3}}\quad\text{(no tail).}
\]
Tracking the $1/\X{1}{3}$ coefficient of each monomial, we find
\[
\bigl(B|_{\Z_1}\bigr)_{1/\X{1}{3}}
=
\frac{a_{\X{1}{3}\X{1}{4}}-a_{\X{2}{4}\X{1}{4}}-a_{\X{2}{5}\X{1}{4}}}{\X{1}{4}}
+\frac{a_{\X{1}{3}\X{3}{5}}-a_{\X{2}{4}\X{3}{5}}-a_{\X{2}{5}\X{3}{5}}}{\X{3}{5}}.
\]
Setting this to zero and using linear independence of $1/\X{1}{4}$ and
$1/\X{3}{5}$:
\begin{align*}
a_{\X{1}{3}\X{1}{4}}&=a_{\X{2}{4}\X{1}{4}}+a_{\X{2}{5}\X{1}{4}},\\
a_{\X{1}{3}\X{3}{5}}&=a_{\X{2}{4}\X{3}{5}}+a_{\X{2}{5}\X{3}{5}}.
\end{align*}
These are \emph{multi-term} relations, not single-coefficient kills.

\paragraph{Step 3: Order $\X{1}{3}^{-2}$ (coefficient of $1/\X{1}{3}^2$).}

This is the most subtle step of the proof and we do it in full.
We enumerate every one of the 15 monomials of $B$, compute its contribution
to the $1/\X{1}{3}^2$ coefficient of $B|_{\Z_1}$, and identify the
single-term kill.

Recall the Laurent expansions:
\[
\frac{1}{\X{2}{4}}=-\frac{1}{\X{1}{3}}-\frac{\X{1}{4}}{\X{1}{3}^2}-\frac{\X{1}{4}^2}{\X{1}{3}^3}-\cdots,\qquad
\frac{1}{\X{2}{5}}=-\frac{1}{\X{1}{3}}\quad\text{(no tail)}.
\]
For products of substituted chords we will also need:
\[
\frac{1}{\X{2}{4}^2}=\frac{1}{\X{1}{3}^2}\cdot\frac{1}{(1-\X{1}{4}/\X{1}{3})^2}=\frac{1}{\X{1}{3}^2}+\frac{2\X{1}{4}}{\X{1}{3}^3}+\cdots,
\]
\[
\frac{1}{\X{2}{4}\X{2}{5}}=-\frac{1}{\X{1}{3}(\X{1}{4}-\X{1}{3})}=\frac{1}{\X{1}{3}^2}+\frac{\X{1}{4}}{\X{1}{3}^3}+\cdots.
\]

The following table lists all 15 monomials of $B$, with the contribution of
each to the coefficient of $1/\X{1}{3}^2$ in $B|_{\Z_1}$:

\begin{center}
\renewcommand{\arraystretch}{1.25}
\begin{tabular}{c|l|l}
\toprule
\# & Monomial of $B$ & Contribution to $(B|_{\Z_1})_{1/\X{1}{3}^2}$ \\
\midrule
1 & $a_{\X{1}{3}\X{1}{3}}/\X{1}{3}^2$ & $+a_{\X{1}{3}\X{1}{3}}$ \\
2 & $a_{\X{1}{3}\X{1}{4}}/(\X{1}{3}\X{1}{4})$ & $0$ (contributes at order $1/\X{1}{3}$) \\
3 & $a_{\X{1}{3}\X{2}{4}}/(\X{1}{3}\X{2}{4})$ & $-a_{\X{1}{3}\X{2}{4}}$ \\
4 & $a_{\X{1}{3}\X{2}{5}}/(\X{1}{3}\X{2}{5})$ & $-a_{\X{1}{3}\X{2}{5}}$ \\
5 & $a_{\X{1}{3}\X{3}{5}}/(\X{1}{3}\X{3}{5})$ & $0$ (order $1/\X{1}{3}$) \\
6 & $a_{\X{1}{4}\X{1}{4}}/\X{1}{4}^2$ & $0$ (order $\X{1}{3}^0$) \\
7 & $a_{\X{1}{4}\X{2}{4}}/(\X{1}{4}\X{2}{4})$ & $-a_{\X{1}{4}\X{2}{4}}$ \\
8 & $a_{\X{1}{4}\X{2}{5}}/(\X{1}{4}\X{2}{5})$ & $0$ (order $1/\X{1}{3}$) \\
9 & $a_{\X{1}{4}\X{3}{5}}/(\X{1}{4}\X{3}{5})$ & $0$ (order $\X{1}{3}^0$) \\
10 & $a_{\X{2}{4}\X{2}{4}}/\X{2}{4}^2$ & $+a_{\X{2}{4}\X{2}{4}}$ \\
11 & $a_{\X{2}{4}\X{2}{5}}/(\X{2}{4}\X{2}{5})$ & $+a_{\X{2}{4}\X{2}{5}}$ \\
12 & $a_{\X{2}{4}\X{3}{5}}/(\X{2}{4}\X{3}{5})$ & $\mathbf{-a_{\X{2}{4}\X{3}{5}}\cdot \X{1}{4}/\X{3}{5}}$ \\
13 & $a_{\X{2}{5}\X{2}{5}}/\X{2}{5}^2$ & $+a_{\X{2}{5}\X{2}{5}}$ \\
14 & $a_{\X{2}{5}\X{3}{5}}/(\X{2}{5}\X{3}{5})$ & $0$ (order $1/\X{1}{3}$) \\
15 & $a_{\X{3}{5}\X{3}{5}}/\X{3}{5}^2$ & $0$ (order $\X{1}{3}^0$) \\
\bottomrule
\end{tabular}
\end{center}

\emph{Verification of a few representative rows.} For row~7,
$1/(\X{1}{4}\X{2}{4}) = (1/\X{1}{4})\cdot(-1/\X{1}{3}-\X{1}{4}/\X{1}{3}^2-\cdots)$;
the $1/\X{1}{3}^2$ piece is $-1/\X{1}{3}^2$, giving coefficient $-1$.
For row~10, the formula above gives $1/\X{2}{4}^2=1/\X{1}{3}^2+O(\X{1}{3}^{-3})$,
so coefficient is $+1$. For row~11 the product identity above gives
$1/(\X{2}{4}\X{2}{5})=+1/\X{1}{3}^2+O(\X{1}{3}^{-3})$, coefficient $+1$.
For row~12,
$1/(\X{2}{4}\X{3}{5})=(1/\X{3}{5})\cdot(-1/\X{1}{3}-\X{1}{4}/\X{1}{3}^2-\cdots)$;
the $1/\X{1}{3}^2$ piece is $-\X{1}{4}/(\X{3}{5}\X{1}{3}^2)$, giving
coefficient $-\X{1}{4}/\X{3}{5}$. Row~13 gives $1/\X{2}{5}^2=1/\X{1}{3}^2$
exactly, coefficient $+1$.

\emph{Summing column 3:}
\[
\bigl(B|_{\Z_1}\bigr)_{1/\X{1}{3}^2}
= \underbrace{\Bigl[a_{\X{1}{3}\X{1}{3}} - a_{\X{1}{3}\X{2}{4}} - a_{\X{1}{3}\X{2}{5}} - a_{\X{1}{4}\X{2}{4}} + a_{\X{2}{4}\X{2}{4}} + a_{\X{2}{4}\X{2}{5}} + a_{\X{2}{5}\X{2}{5}}\Bigr]}_{K\ :=\ \text{constant}\ (\text{no free-var dependence})} - a_{\X{2}{4}\X{3}{5}}\cdot\frac{\X{1}{4}}{\X{3}{5}}.
\]

This entire expression must vanish identically as a rational function of
two independent variables $(\X{1}{4},\X{3}{5})$. The two rational functions
\[
1 \quad\text{and}\quad \frac{\X{1}{4}}{\X{3}{5}}
\]
are linearly independent: the first is constant, the second grows linearly
in $\X{1}{4}$ and has a simple pole at $\X{3}{5}=0$. Therefore each
coefficient must vanish independently:
\begin{itemize}[nosep,leftmargin=2em]
\item The coefficient of $1$: $\quad K = 0$ --- a multi-term equation
relating $7$ coefficients.
\item The coefficient of $\X{1}{4}/\X{3}{5}$:
\[
\boxed{a_{\X{2}{4}\X{3}{5}}=0\qquad\text{($\Z_1$, single-term order-2 kill).}}
\]
\end{itemize}

\paragraph{The mechanism, named.} Observe the critical structural feature:
only one monomial of $B$ (row~12) produced a non-constant free-variable
rational function at order $1/\X{1}{3}^2$. This is no accident. It happens
because:
\begin{enumerate}[nosep,leftmargin=2em]
\item The non-bare substitution $\X{2}{4}=\X{1}{4}-\X{1}{3}$ has a
\textbf{companion} $\X{1}{4}$ appearing at the $1/\X{1}{3}^2$ level of
its Laurent tail.
\item For a monomial $a_{\X{2}{4}X_b}/(\X{2}{4}X_b)$ with $X_b$ free, this
tail contributes the factor $-\X{1}{4}/X_b$ --- a non-constant rational
function of $X_b$ if $X_b\ne \X{1}{4}$.
\item If $X_b = \X{1}{4}$ itself (the companion), the contribution degenerates
to a constant $-1$, and this source loses its distinctness against
all the other constant contributions (rows 1, 3, 4, 7, 10, 11, 13).
\end{enumerate}

This gives us a general rule:

\begin{lemma}[Companion-isolation rule, single-companion case]
\label{lem:companion-isolation-simple}
Let $\Z_r$ have a non-bare substituted chord $X_s = X_f - X_*$ with
companion $X_f \in F_r$. For each free chord $X_b \in F_r \setminus \{X_f\}$,
the monomial $X_f/X_b$ is a unique source at order $1/X_*^{2}$ in the
Laurent expansion of $B|_{\Z_r}$. Consequently,
\[
a_{X_s\,X_b}\ =\ 0\quad\text{for every}\quad X_b \in F_r\setminus\{X_f\}.
\]
\end{lemma}

At $n=5$ with $|F_r|=2$ and one non-bare substitution per zone, the lemma
produces $|F_r|-1=1$ kill per zone, accounting for all five order-2 kills
in the cyclic rotation.

\medskip

\begin{remark}[Preview of generalization]
At $n=6$ each zone has \emph{two} non-bare substitutions, each with its own
companion. The companion-isolation lemma applies independently to each,
giving two separate families of $|F_r|-1$ kills at order 2.
Moreover, a new layer appears at order 4, coming from monomials with
\emph{both} non-bare indices simultaneously; this will require a
\emph{double}-companion isolation argument using the product Laurent
expansion $1/(\text{non-bare}_1\cdot\text{non-bare}_2)$. We set this up
after completing the $n=5$ proof.
\end{remark}

\paragraph{Step 4: orders $\X{1}{3}^{-3}, \X{1}{3}^{-4}, \ldots$.}

A direct check shows that no new single-coefficient kills arise. The same
monomial $\X{1}{4}/\X{3}{5}$ recurs (giving $a_{\X{2}{4}\X{3}{5}}=0$
redundantly), and new monomials produce multi-term equations.

\paragraph{Summary of $\Z_1$:}
\[
\text{Single-coefficient kills from }\Z_1:\quad
\boxed{a_{\X{1}{4}\X{1}{4}},\ a_{\X{3}{5}\X{3}{5}},\ a_{\X{1}{4}\X{3}{5}}}\ \text{(order 0)},\quad
\boxed{a_{\X{2}{4}\X{3}{5}}}\ \text{(order 2)}.
\]
Four kills total.

\subsection{Cycling: kills from \texorpdfstring{$\Z_2,\ldots,\Z_5$}{Z2..Z5}}

By cyclic symmetry, each $\Z_r$ has the same structure as $\Z_1$, rotated.
Reading off:

\begin{center}
\renewcommand{\arraystretch}{1.2}
\begin{tabular}{c|c|c}
\toprule
Zero & Order-0 kills (multisets in $F_r$) & Order-2 kill \\
\midrule
$\Z_1$ ($X_*=\X{1}{3}$) & $a_{\X{1}{4}\X{1}{4}},\ a_{\X{3}{5}\X{3}{5}},\ a_{\X{1}{4}\X{3}{5}}$ & $a_{\X{2}{4}\X{3}{5}}$\\
$\Z_2$ ($X_*=\X{2}{4}$) & $a_{\X{1}{4}\X{1}{4}},\ a_{\X{2}{5}\X{2}{5}},\ a_{\X{1}{4}\X{2}{5}}$ & $a_{\X{1}{4}\X{3}{5}}$\\
$\Z_3$ ($X_*=\X{3}{5}$) & $a_{\X{1}{3}\X{1}{3}},\ a_{\X{2}{5}\X{2}{5}},\ a_{\X{1}{3}\X{2}{5}}$ & $a_{\X{1}{4}\X{2}{5}}$\\
$\Z_4$ ($X_*=\X{1}{4}$) & $a_{\X{1}{3}\X{1}{3}},\ a_{\X{2}{4}\X{2}{4}},\ a_{\X{1}{3}\X{2}{4}}$ & $a_{\X{1}{3}\X{2}{5}}$\\
$\Z_5$ ($X_*=\X{2}{5}$) & $a_{\X{2}{4}\X{2}{4}},\ a_{\X{3}{5}\X{3}{5}},\ a_{\X{2}{4}\X{3}{5}}$ & $a_{\X{1}{3}\X{2}{4}}$\\
\bottomrule
\end{tabular}
\end{center}

\subsection{Union: every non-local coefficient is killed}

\paragraph{Double poles (5 total):} each is killed at order 0 by exactly two zones.
\begin{itemize}[nosep,leftmargin=2em]
\item $a_{\X{1}{3}^2}$: by $\Z_3$ and $\Z_4$.
\item $a_{\X{1}{4}^2}$: by $\Z_1$ and $\Z_2$.
\item $a_{\X{2}{4}^2}$: by $\Z_4$ and $\Z_5$.
\item $a_{\X{2}{5}^2}$: by $\Z_2$ and $\Z_3$.
\item $a_{\X{3}{5}^2}$: by $\Z_1$ and $\Z_5$.
\end{itemize}

\paragraph{Crossing pairs (5 total):} each is killed once at order 0 and once at order 2.
\begin{itemize}[nosep,leftmargin=2em]
\item $a_{\X{1}{3}\X{2}{4}}$: order 0 by $\Z_4$, order 2 by $\Z_5$.
\item $a_{\X{1}{3}\X{2}{5}}$: order 0 by $\Z_3$, order 2 by $\Z_4$.
\item $a_{\X{1}{4}\X{2}{5}}$: order 0 by $\Z_2$, order 2 by $\Z_3$.
\item $a_{\X{1}{4}\X{3}{5}}$: order 0 by $\Z_1$, order 2 by $\Z_2$.
\item $a_{\X{2}{4}\X{3}{5}}$: order 0 by $\Z_5$, order 2 by $\Z_1$.
\end{itemize}

\paragraph{Tree triples (5 total) are never killed.} For each fan triangulation
$a_{\X{i}{j}\X{j}{k}}$ (chords sharing endpoint $j$), at every zone $\Z_r$
at least one of the two indices is either the special $X_*$, the bare
substituted chord, or one of the substituted chords adjacent to the special
in a way that breaks the conditions for both order-0 and order-2 kills.
Equivalently: the tree amplitude $A_5$ vanishes on each $\Z_r$, so every
Laurent coefficient of $A_5|_{\Z_r}$ vanishes; the tree triples
in $A_5$ form the kernel of all our conditions.

\begin{theorem}[Locality at five points]
The five hidden one-zeros of the pentagon force every non-local coefficient
of the general rational ansatz $B$ to vanish. The surviving 5 coefficients
are precisely the tree triangulations.
\end{theorem}

\begin{proof}
By the explicit calculation at $\Z_1$ above (Steps 1--4) and cyclic rotation
to obtain $\Z_2,\ldots,\Z_5$, the union of single-coefficient kills covers
all 10 non-local coefficients (5 double poles, 5 crossing pairs), as
tabulated. The 5 tree triples are never killed. \qed
\end{proof}

\subsection{The geometric pattern revealed}
\label{sec:geometric-pattern-n5}

Look at the table of which zone kills each crossing pair at which order.
A clean pattern emerges: every crossing pair is killed by exactly two zones,
and these zones are \emph{adjacent} in the cyclic ordering. The order-0
kill comes from the zone where both chords lie in the free set; the order-2
kill comes from the next zone over, where the special variable becomes one
of the chord-pair's neighbors via the enclosed-vertex correspondence.

This is the \textbf{enclosed-vertex correspondence in action}: the hidden
zeros are tightly linked to the crossings of the polygon, and every
crossing has a precise pair of killing zones determined by the
enclosed-vertex structure.

%================================================================
\section{Six points: the hexagon case}
\label{sec:n6}

The same technique works at $n=6$, with two new features:
\begin{itemize}[nosep,leftmargin=2em]
\item The hexagon has \emph{two classes} of chords (short and long diagonals).
\item Each zero now has \emph{three} substitutions, of which \emph{two} are
non-bare. This produces an additional ``order 4'' layer of kills coming
from monomials with two non-bare substituted indices.
\end{itemize}

\subsection{Setup at $n=6$}

The hexagon's 9 chords:
\begin{center}
\begin{tabular}{ll}
\textbf{Short diagonals} (class 1, enclose 1 vertex): &
$\X{1}{3},\X{2}{4},\X{3}{5},\X{4}{6},\X{1}{5},\X{2}{6}$ \\
\textbf{Long diagonals} (class 2, enclose 2 vertices): &
$\X{1}{4},\X{2}{5},\X{3}{6}$
\end{tabular}
\end{center}

Cyclic shorthand:
\[
x_1=\X{1}{3},\ x_2=\X{2}{4},\ x_3=\X{3}{5},\ x_4=\X{4}{6},\ x_5=\X{1}{5},\ x_6=\X{2}{6};
\]
\[
y_1=\X{1}{4},\ y_2=\X{2}{5},\ y_3=\X{3}{6}.
\]

Enclosed-vertex sets:
\[
E(x_r) = \{r+1\}\bmod 6,\qquad E(y_r) = \{r+1, r+2\}\bmod 6.
\]

\paragraph{Ansatz.} Each monomial has $n-3=3$ chord factors. The number of
multisets of size 3 from 9 chords is $\binom{9+2}{3}=165$. Of these, $14$
are tree triples ($C_4=14$ triangulations of the hexagon), $151$ are
non-local.

\subsection{The six zeros, derived by hand}

Each $\Z_r$ is defined by $c_{r,r+2}=c_{r,r+3}=c_{r,r+4}=0$ (three constraints).
For $\Z_1$:

$c_{1,3}=\X{1}{4}+\X{2}{3}-\X{1}{3}-\X{2}{4}=y_1-x_1-x_2=0\Rightarrow x_2=y_1-x_1$.

$c_{1,4}=\X{1}{5}+\X{2}{4}-\X{1}{4}-\X{2}{5}=x_5+x_2-y_1-y_2=0$.
Using $x_2=y_1-x_1$: $x_5+y_1-x_1-y_1-y_2=x_5-x_1-y_2=0$, so $y_2=x_5-x_1$.

$c_{1,5}=\X{1}{6}+\X{2}{5}-\X{1}{5}-\X{2}{6}=0+y_2-x_5-x_6=0$.
Using $y_2=x_5-x_1$: $-x_1-x_6=0$, so $x_6=-x_1$.

So $\Z_1$: $x_2=y_1-x_1$, $y_2=x_5-x_1$, $x_6=-x_1$, with special $x_*=x_1=\X{1}{3}$.
The substituted chords $\{x_2,y_2,x_6\}=\{\X{2}{4},\X{2}{5},\X{2}{6}\}$ are
exactly the chords from the enclosed vertex $2$ of $\X{1}{3}$.

By identical derivations (cyclic), the six zeros are:

\begin{center}
\renewcommand{\arraystretch}{1.2}
\begin{tabular}{c|c|c|c}
\toprule
Zero & Special & Substitutions & Free $F_r$ (excluding spec.)\\
\midrule
$\Z_1$ & $x_1$ & $x_2=y_1-x_1$, $y_2=x_5-x_1$, $x_6=-x_1$ & $\{x_3,x_4,x_5,y_1,y_3\}$\\
$\Z_2$ & $x_2$ & $x_3=y_2-x_2$, $y_3=x_6-x_2$, $x_1=-x_2$ & $\{x_4,x_5,x_6,y_1,y_2\}$\\
$\Z_3$ & $x_3$ & $x_4=y_3-x_3$, $y_1=x_1-x_3$, $x_2=-x_3$ & $\{x_1,x_5,x_6,y_2,y_3\}$\\
$\Z_4$ & $x_4$ & $x_5=y_1-x_4$, $y_2=x_2-x_4$, $x_3=-x_4$ & $\{x_1,x_2,x_6,y_1,y_3\}$\\
$\Z_5$ & $x_5$ & $x_6=y_2-x_5$, $y_3=x_3-x_5$, $x_4=-x_5$ & $\{x_1,x_2,x_3,y_1,y_2\}$\\
$\Z_6$ & $x_6$ & $x_1=y_3-x_6$, $y_1=x_4-x_6$, $x_5=-x_6$ & $\{x_2,x_3,x_4,y_2,y_3\}$\\
\bottomrule
\end{tabular}
\end{center}

In every zone:
\begin{itemize}[nosep,leftmargin=2em]
\item One special $x_*$ (a short diagonal).
\item Three substituted chords from the enclosed vertex of $x_*$. Of these:
  \emph{two are non-bare} (one short, one long, with companions in $F_r$),
  and \emph{one is bare} ($=-x_*$).
\item Five free chords $F_r$ (three short, two long).
\end{itemize}

\subsection{The kill technique at \texorpdfstring{$\Z_1$}{Z1}: three layers}

For $\Z_1$: $x_*=x_1$, non-bare substitutions
\[
x_2=y_1-x_1\ \text{(companion $f_{x_2}=y_1$)},\quad y_2=x_5-x_1\ \text{(companion $f_{y_2}=x_5$)},
\]
bare: $x_6=-x_1$. Free: $F_1=\{x_3,x_4,x_5,y_1,y_3\}$ (5 chords).

The Laurent expansions:
\[
\frac{1}{x_2}=-\frac{1}{x_1}-\frac{y_1}{x_1^2}-\frac{y_1^2}{x_1^3}-\cdots,
\quad
\frac{1}{y_2}=-\frac{1}{x_1}-\frac{x_5}{x_1^2}-\frac{x_5^2}{x_1^3}-\cdots,
\quad
\frac{1}{x_6}=-\frac{1}{x_1}.
\]

\paragraph{Layer 1: Order $x_1^0$ ($x_1\to\infty$ limit). [\textbf{35 kills}]}

As $x_1\to\infty$, all of $x_2,y_2,x_6$ tend to $\pm\infty$, so $1/x_2,1/y_2,1/x_6\to 0$.
Surviving monomials of $B$ are those with all three indices in $F_1$:
\[
\bigl(B|_{\Z_1}\bigr)\bigr|_{x_1\to\infty} = \sum_{(a,b,c)\subseteq F_1}\frac{a_{abc}}{X_a X_b X_c}.
\]
The number of multisets $(a,b,c)\subseteq F_1$ (size 3 from 5 elements) is $\binom{5+2}{3}=35$.
Each $1/(X_aX_bX_c)$ is a distinct rational monomial in 5 free variables;
they are linearly independent. Hence each $a_{abc}=0$:
\[
\boxed{\text{35 single-term kills at order 0: all }a_{abc}\text{ with }\{X_a,X_b,X_c\}\subseteq F_1.}
\]

\paragraph{Layer 2: Order $x_1^{-2}$ (coefficient of $1/x_1^2$). [\textbf{20 new kills}]}

This is where the companion-isolation rule (Lemma~\ref{lem:companion-isolation-simple})
applies \emph{twice}, once for each non-bare substituted chord.

Recall the statement: for a non-bare substitution $X_s = X_f - X_*$ with
companion $X_f$, the monomial $X_f/X_b$ at order $1/X_*^2$ has a unique
source among the terms of $B$, provided $X_b \ne X_f$. Pairing $X_s$ with
any two free chords $X_a, X_b \in F_r \setminus \{X_f\}$ gives a unique
monomial $X_f/(X_a X_b)$ at order $1/x_1^2$.

\paragraph{Generalized companion-isolation (two-factor version).} The only
subtlety at $n=6$ is that now we have \emph{three} factors in each monomial,
but only one of them is the substituted $X_s$; the other two are the
``bystanders'' in $F_r \setminus \{X_f\}$. The Laurent-tail contribution
still produces the unique rational monomial $X_f/(X_a X_b)$ at order
$1/x_1^2$, just with two free denominators instead of one.

\begin{lemma}[Companion-isolation, general form]
\label{lem:companion-isolation-general}
Let $\Z_r$ have a non-bare substituted chord $X_s = X_f - X_*$. For each
multiset $(X_a, X_b, \ldots)$ of size $(n-3)-1$ drawn from $F_r \setminus \{X_f\}$,
the monomial $X_f/(X_a X_b\cdots)$ at order $1/X_*^2$ in the Laurent
expansion of $B|_{\Z_r}$ has a unique source, namely
$a_{X_s X_a X_b \cdots}/(X_s X_a X_b\cdots)$. Consequently, the kill is:
\[
a_{X_s\,X_a X_b\cdots}\ =\ 0
\quad\text{for every multiset in }F_r\setminus\{X_f\}\text{ of size }n-4.
\]
\end{lemma}

\emph{Proof (sketch).} The Laurent expansion of $1/X_s=1/(X_f-X_*)$ at
order $1/X_*^2$ contributes $-X_f/X_*^2$. Multiplying by the remaining
$1/(X_a X_b\cdots)$ gives $-X_f/(X_*^2 X_a X_b\cdots)$. The monomial
$X_f/(X_a X_b\cdots)$ is non-degenerate iff no $X_a, X_b, \ldots$ equals
$X_f$ (else the $X_f$ ``cancels'' and we get a constant monomial,
re-entering the multi-term pool). For any other monomial of $B$ to
contribute this exact rational function would require:
\begin{itemize}[nosep,leftmargin=2em]
\item Either a different non-bare substitution $X_{s'}=X_{f'}-X_*$ with
$f'=f$ --- impossible, since each non-bare has a distinct companion;
\item Or a combination of Laurent tails and free factors producing
$X_f/(X_a X_b \cdots)$ --- these are higher-order tails or repeated-$X_s$
contributions that only enter at orders $\ge 3$, outside the $1/X_*^2$ piece.
\end{itemize}
\qed

Applying Lemma~\ref{lem:companion-isolation-general} to $\Z_1$ at $n=6$:

\textbf{Sub-case A: $X_s = x_2$ (non-bare, companion $f_{x_2}=y_1$).}
Pair with multisets of size $(n-3)-1 = 2$ from $F_1\setminus\{y_1\}=\{x_3,x_4,x_5,y_3\}$
($4$ elements). Number of multisets of size 2 from 4 elements:
$\binom{4+1}{2}=10$. These ten multisets give
\[
\boxed{10\text{ single-term kills: }a_{x_2 X_a X_b} = 0 \text{ for } \{X_a,X_b\}\subseteq\{x_3,x_4,x_5,y_3\}.}
\]

\textbf{Sub-case B: $X_s = y_2$ (non-bare, companion $f_{y_2}=x_5$).}
Multisets of size 2 from $F_1\setminus\{x_5\}=\{x_3,x_4,y_1,y_3\}$:
another $\binom{4+1}{2}=10$ kills.
\[
\boxed{10\text{ more single-term kills: }a_{y_2 X_a X_b} = 0 \text{ for } \{X_a,X_b\}\subseteq\{x_3,x_4,y_1,y_3\}.}
\]

\textbf{Sub-case C: $X_s=x_6$ (bare).} Since $1/x_6=-1/x_1$ has no Laurent
tail beyond leading order, it contributes $0$ to the order-$1/x_1^2$
companion structure. No kills from this sub-case.

\textbf{Total at order 2: $10+10=20$ kills.}

\paragraph{Layer 3: Order $x_1^{-4}$ (coefficient of $1/x_1^4$). [\textbf{3 new kills}]}

This is the new layer that appears at $n=6$ but was absent at $n=5$ (where
there was only one non-bare substitution per zone). It uses both non-bare
substitutions simultaneously --- a \emph{double companion-isolation}.

\paragraph{Double companion-isolation.} Consider monomials of $B$ with
\emph{both} non-bare substitutions as factors: $a_{x_2\,y_2\,X_c}/(x_2\,y_2\,X_c)$
with $X_c\in F_1$.

Use the product Laurent expansion:
\[
\frac{1}{x_2\,y_2}=\frac{1}{(y_1-x_1)(x_5-x_1)}=\frac{1}{x_1^2}\cdot\frac{1}{(1-y_1/x_1)(1-x_5/x_1)}.
\]
Expanding each factor and multiplying:
\[
\frac{1}{(1-y_1/x_1)(1-x_5/x_1)}
= 1 + \frac{y_1+x_5}{x_1} + \frac{y_1^2+y_1 x_5+x_5^2}{x_1^2} + O(x_1^{-3}).
\]
So
\[
\frac{1}{x_2\,y_2}\bigg|_{1/x_1^4} = \frac{y_1^2+y_1 x_5+x_5^2}{x_1^4}.
\]
The coefficient at $1/x_1^4$ from $a_{x_2\,y_2\,X_c}/(x_2\,y_2\,X_c)$ is therefore
\[
\frac{a_{x_2\,y_2\,X_c}\,(y_1^2+y_1 x_5+x_5^2)}{X_c}.
\]

\paragraph{The key isolating monomial: $y_1\,x_5/X_c$.} Of the three
``free-variable'' monomials produced at this order --- $y_1^2/X_c$,
$y_1 x_5/X_c$, $x_5^2/X_c$ --- the \emph{mixed} one $y_1\,x_5/X_c$
is the uniquely-identifying fingerprint:
\begin{itemize}[nosep,leftmargin=2em]
\item $y_1^2/X_c$ can also arise from $1/x_2^2$-type monomials at order $1/x_1^4$
(which contribute pure $y_1^k$ monomials from the Laurent tail of $1/x_2^2$).
\item $x_5^2/X_c$ similarly arises from $1/y_2^2$-type monomials.
\item But $y_1\,x_5/X_c$ \textbf{requires} one factor of the $x_2$-tail
and one factor of the $y_2$-tail --- i.e., the monomial of $B$ \emph{must}
contain both $x_2$ and $y_2$ as factors. No other source produces this
mixed monomial at order $1/x_1^4$.
\end{itemize}

\emph{Non-degeneracy condition.} As before, the monomial $y_1\,x_5/X_c$ is
non-degenerate (a genuine rational function, not a constant) iff
$X_c \notin \{y_1, x_5\}$. That is, $X_c\in F_1\setminus\{y_1,x_5\}=\{x_3,x_4,y_3\}$.
So the kills are:
\[
\boxed{\text{3 single-term kills at order 4: }
a_{x_2\,y_2\,x_3}=a_{x_2\,y_2\,x_4}=a_{x_2\,y_2\,y_3}=0.}
\]

This is the double companion-isolation. The mixed monomial
$f_{s_1}\cdot f_{s_2}/X_c$ (product of the two companions) is the fingerprint
that identifies monomials carrying both non-bare substitutions at order
$1/X_*^4$.

\paragraph{Higher orders contribute no new single-term kills.}

The monomials with three non-bare/bare-factor combinations either
involve $x_6$ (bare, no tail) or repeat the same Laurent structure already
captured. Direct expansion shows no new isolable single-term equations.

\paragraph{Total kills from $\Z_1$: $35+20+3 = 58$ single-coefficient kills.}

\subsection{Cycling: union over six zones is 151}

By cyclic symmetry every $\Z_r$ produces $35+20+3=58$ kills, with the same
internal structure (just rotated). Tabulating all $6\times 58=348$ raw
kills and accounting for overlaps yields a union of $151$ distinct
coefficients --- exactly the number of non-local multisets at $n=6$.

\paragraph{Sketch of the union argument.} It suffices to show every
non-local multiset is killed by at least one zone. The non-local multisets
at $n=6$ split into types by chord-class composition (\#short, \#long):

\begin{center}
\begin{tabular}{l|c|c}
\toprule
Composition & \# multisets & \# trees among them \\
\midrule
3 short & $\binom{6+2}{3}=56$ & 2 (the two ``zigzag'' triangulations $x_1x_3x_5$, $x_2x_4x_6$) \\
2 short, 1 long & $3\cdot \binom{6+1}{2}=63$ & 6 (six fan-like) \\
1 short, 2 long & $6\cdot \binom{3+1}{2}=36$ & 6 \\
3 long & $\binom{3+2}{3}=10$ & 0 \\
\midrule
Total & 165 & 14 \\
\bottomrule
\end{tabular}
\end{center}

For each non-tree multiset, one identifies a zone $\Z_r$ where the
multiset's indices fit one of the three layers (all in $F_r$;
exactly one non-bare $s$ with the rest in $F_r\setminus\{f_s\}$; or both
non-bare $s$'s with the third in $F_r\setminus\{f_{s_1},f_{s_2}\}$).
A full case analysis confirms every non-tree multiset is caught.

The 14 tree triples are not killed: each is a triangulation (non-crossing
chords), so for every zone the multiset includes either the special $X_*$
itself, or a chord that violates the unique-source conditions of the
three layers.

\begin{theorem}[Locality at six points]
The six hidden one-zeros of the hexagon force every non-local coefficient
of the general rational ansatz $B$ to vanish. The $14$ surviving
coefficients are exactly the tree triangulations of the hexagon.
\end{theorem}

%================================================================
\section{Looking ahead: the structural pattern}

The kill structure at $n=4,5,6$ reveals a clean iterated mechanism.
At each zone $\Z_r$, there are:
\begin{itemize}[nosep,leftmargin=2em]
\item $|F_r|=\frac{n(n-3)}{2}-(n-3)-1$ free chords (the non-substituted,
non-special ones).
\item $n-4$ \emph{non-bare} substituted chords, each with its own companion
in $F_r$ (at $n=4$ there are $0$ non-bares; at $n=5$ there is $1$;
at $n=6$ there are $2$; and in general at $n$ points, $n-4$).
\item $1$ bare substituted chord (contributes only to the leading Laurent order).
\end{itemize}

The single-term kills stratify into $k$-fold companion-isolation layers:
\begin{itemize}[nosep,leftmargin=2em]
\item \textbf{Layer $0$ (order $X_*^0$):} the $X_*\to\infty$ limit kills
all multisets with indices entirely in $F_r$.
Count: $\binom{|F_r|+n-4}{n-3}$.
\item \textbf{Layer $1$ (order $1/X_*^2$):} single companion-isolation
(Lemma~\ref{lem:companion-isolation-general}) kills multisets with
exactly one non-bare index $X_s$ and the remaining indices in
$F_r\setminus\{f_s\}$. Count: $(n-4)\cdot\binom{|F_r|+n-6}{n-4}$.
\item \textbf{Layer $k$ (order $1/X_*^{2k}$):} $k$-fold companion-isolation
kills multisets with exactly $k$ non-bare indices, using the product
Laurent expansion
$1/(X_{s_1}\cdots X_{s_k}) = (1/X_*^k)\cdot\prod_i(1-f_{s_i}/X_*)^{-1}$
whose $1/X_*^{2k}$ piece produces the mixed monomial
$f_{s_1}f_{s_2}\cdots f_{s_k}/\prod X_b$ as the unique source of the kill.
Count: $\binom{n-4}{k}\cdot\binom{|F_r|-k+(n-3-k)-1}{n-3-k}$.
\end{itemize}

The per-zone total is the sum:
\[
\text{kills per zone}\;=\;\sum_{k=0}^{\lfloor (n-3)/2\rfloor}
\binom{n-4}{k}\cdot\binom{|F_r|-k+(n-3-k)-1}{n-3-k}.
\]
At $n=4,5,6$ this gives $0,\ 4,\ 58$ kills per zone, with respective unions
$0,\ 10,\ 151$ --- exactly the number of non-local coefficients.

At $n=7$ the formula predicts
$495+360+84+6=945$ kills per zone, coming from four layers (up to $k=3$
since there are $n-4=3$ non-bare substitutions). The union size (and the
statement of locality) we leave as the next target.

This structural template suggests a natural induction on $n$: both the
enclosed-vertex principle and the $k$-fold companion-isolation mechanism
scale cleanly as $n$ grows. The class hierarchy of chords (short, long,
extra-long, etc.) --- each new class being born at $n=2k$ --- refines the
bookkeeping further.

\end{document}
